
This paper investigates whether machine learning benchmark score compression — the narrowing of discriminative score distributions on leaderboards — can be detected systematically and whether submission accumulation temporally precedes such compression. Using a quarterly panel of 6,938 benchmark-quarter observations across 466 qualifying benchmarks from the Papers With Code HuggingFace archive (2018–2025), 31.1\% of qualifying benchmarks (145 of 466) exhibit at least one compression event, defined as score variance among the top-10 models falling below 1.5 times the estimated measurement noise for two or more consecutive quarters. Granger causality analysis indicates that cumulative submission count temporally precedes score variance compression at lag 2 (approximately 6 months; p = 1.854 × 10⁻⁵), while the reverse direction is not confirmed. Domain-specific health estimators (H_d) were evaluated on synthetic benchmark panels calibrated to real data statistics; Mann-Whitney U tests yielded p < 0.0001 with |Cohen's d| > 5 in all three domains (CV, NLP, tabular), though these effect sizes have not been validated on real benchmark data and may be optimistic due to the idealized nature of the synthetic panels. Cross-sectional analysis on real data showed that NLP H_d achieves an AUC of 0.857 for distinguishing compressed from non-compressed benchmarks. The planned Cox proportional hazards survival model (H-M3) was not executed due to a collapse event operationalization failure: only 1 event was detected under the pre-registered expanding Kendall τ criterion (far below the required minimum of 20), blocking H-M3 and H-M4 execution. The prospective prediction component of the proposed framework — including C-index ≥ 0.70 and early warning lead time ≥ 12 months — remains empirically unvalidated in this work.

